\documentclass[11pt]{report}
\input{../pri-end/T0_preamble_standalone_A4_De}
% T0_preamble_patches.tex — LOKALER PLATZHALTER
% Im Repo durch die offizielle Version ersetzen.


\title{Dok.~313 -- Kein Anfang:\\
Die thermische Geschichte auf dem Zeitzyklus}
\author{Johann Pascher}
\date{August 2026}

\begin{document}
\maketitle
\tableofcontents

% ============================================================
\chapter*{Dok.~313 -- Kein Anfang}
\addcontentsline{toc}{chapter}{Dok.~313 -- Kein Anfang}
% ============================================================

\section*{Worum es geht}

Dok.~312 endet mit Pflicht~D(ii): Können die Feldkonfigurationen
auf dem kompakten Zeitzyklus ($\tau_c = 2\pi/H_0 = 91{,}9$~Gyr)
periodisch schließen \emph{und dabei} die beobachtete thermische
Abfolge tragen -- Nukleosynthese, Rekombination, Strukturbildung?
Ein Kreis hat keinen Anfang; die Beobachtungen zeigen aber einen
Gradienten. Dieses Dokument rechnet die Frage durch
(\texttt{ffgft\_313\_zeitzyklus\_geschichte.py}) und findet drei
Ergebnisse und einen harten offenen Kern:

\begin{enumerate}
  \item \textbf{Antipoden-Befund [K$|$P20]:} Die gesamte
        beobachtbare Geschichte füllt auf $1$--$2\,\%$ genau
        \emph{einen Halbzyklus}. Die letzte Streufläche sitzt bei
        $\theta = 0{,}992\,\pi$, der Partikelhorizont bei
        $1{,}012\,\pi$ -- der \enquote{Anfang} ist der Antipode des
        Zeitzyklus, kein Ereignis.
  \item \textbf{Glatte Schließung [K]:} Das beobachtete
        Massenlauf-Profil erreicht den Antipoden mit Steigung
        $\sqrt{\Omega_r} \approx 0{,}01$ -- die
        Glattheitsbedingung einer geraden periodischen Schließung
        ist auf $1\,\%$ erfüllt, und die Rückkehr-Hälfte ist exakt
        die photonisch unbeobachtbare Hälfte.
  \item \textbf{Entropie-Hindernis [K]:} Exakte
        $\tau_c$-Periodizität ist dynamisch nicht generisch
        (Poincaré-Diskrepanz $\sim 10^{10^{104}}$); sie muss als
        Randbedingung gesetzt werden. Feinkörnig zulässig,
        grobkörnig offen -- das ist der präzisierte Kern von
        D(ii).
\end{enumerate}

Kein Teil der A-Serie; Registrierung in Dok.~190 separat. Bezüge:
Dok.~306 ($\partial\text{T}^4 = \emptyset$), 267 (Entartung),
312 (Abschlussskala, kompakte Zeitrichtung, Statik der Relation,
Torus-Ruhesystem).

% ============================================================
\chapter{Kein Rand, kein $t=0$: die Ausgangslage}
% ============================================================

Die Position der FFGFT zu Singularität und Urknall steht seit
Dok.~306: $\partial\text{T}^4 = \emptyset$ -- der Torus hat keinen
Rand, \enquote{was war bei $t=0$?} ist falsch gestellt. Dok.~312
hat diese Position konkretisiert und dreifach gestützt:

\textbf{Geometrisch:} Die Zeitrichtung ist ein Zyklus der Länge
$91{,}9$~Gyr. Ein geschlossener Kreis hat keinen Anfangspunkt --
\enquote{$t=0$} ist auf ihm so sinnlos wie \enquote{der Anfang des
Äquators}.

\textbf{Rückwärts:} In der Massenlauf-Lesart (Dok.~312, Statik der
Relation; Wetterich) gehen zurückgerechnet nicht Abstände gegen
null, sondern Massen: $\tilde T \cdot m = 1$ heißt dann
$\tilde T \to \infty$ -- keine Dichtedivergenz, sondern ein Regime
immer langsamerer Uhren. Die \enquote{Singularität} ist der Ort,
an dem die Uhren stehen, nicht an dem die Dichte divergiert.

\textbf{Vorwärts:} Kein Kollaps in eine Singularität -- das
Windungs-/Impuls-Minimum (Anhang Dok.~312) hat $V'' > 0$.

Was Dok.~312 offen ließ, ist die \emph{Geschichte}: Die
Beobachtungen zeigen eine klare thermische Abfolge. Wie verträgt
sich ein Gradient mit einem Kreis? Das ist D(ii), und es wird hier
gerechnet, nicht erzählt.

% ============================================================
\chapter{Der Antipoden-Befund [K$|$P20]}
% ============================================================

\section{Wo sitzt die Geschichte auf dem Zyklus?}

Position auf dem Zyklus: $\theta = \chi/R_H$, mit $\chi$ als
Lichtweg-Distanz (die $z \to \chi$-Zuordnung übernimmt die
gemessene $\Lambda$CDM-Abbildung; nur dimensionslose Verhältnisse
gehen ein, Dok.~267). Der Antipode liegt bei $\theta = \pi$, d.\,h.
$\chi = L^{*}/2 = \pi R_H = 14{,}09$~Gpc. Gerechnet:

\begin{center}
\begin{tabular}{lcc}
\toprule
 & $\chi$ [Gpc] & Anteil von $L^{*}/2$ \\
\midrule
Rekombination ($z=1090$) & $13{,}98$ & $0{,}992$ \\
Partikelhorizont ($z\to\infty$) & $14{,}27$ & $1{,}012$ \\
\bottomrule
\end{tabular}
\end{center}

\begin{center}
\begin{tabular}{lcc}
\toprule
$z$ & $\theta/\pi$ & $m(z)/m_0 = 1/(1{+}z)$ \\
\midrule
$0{,}5$ & $0{,}140$ & $0{,}667$ \\
$1$ & $0{,}243$ & $0{,}500$ \\
$3$ & $0{,}465$ & $0{,}250$ \\
$10$ & $0{,}689$ & $0{,}091$ \\
$100$ & $0{,}917$ & $0{,}0099$ \\
$1090$ & $0{,}992$ & $0{,}00092$ \\
\bottomrule
\end{tabular}
\end{center}

\textbf{Befund:} Die gesamte beobachtbare Geschichte -- bis
$z \to \infty$ -- füllt auf $1{,}2\,\%$ genau einen
\emph{Halbzyklus}. Die letzte Streufläche sitzt am Antipoden
($99{,}2\,\%$), der formale \enquote{Anfang} $1{,}2\,\%$ dahinter.
Der \enquote{Urknall} ist damit ein \emph{Ort auf dem Zyklus} --
die Antipodenregion, in der Massenlauf-Lesart: die Region der
kleinsten Massen und langsamsten Uhren -- und kein Ereignis vor
der Zeit.

\section{Ehrliche Einordnung der Koinzidenz [S]}

In $\Lambda$CDM-Sprache ist derselbe Sachverhalt die bekannte
Quasi-Koinzidenz $\eta_0 H_0/c \approx \pi$ (hier: $3{,}18$ gegen
$3{,}14$, also $1{,}2\,\%$) -- dort eine zufällige Folge der
gemessenen $\Omega$-Werte. In der kompakten Lesart wird daraus
eine Strukturaussage: \emph{Die Streuwand sitzt am Antipoden.}
Damit sie mehr ist als eine umgedeutete Koinzidenz, müsste die
Lage vorwärts folgen -- etwa daraus, dass der Materiegehalt die
Wandposition an den Halbzyklus bindet. Das ist nicht gezeigt;
Status [S], als Herleitungsziel notiert.

% ============================================================
\chapter{Die glatte Schließung [K]}
% ============================================================

\section{Das Profil und seine Endpunktsteigung}

Periodizität verlangt, dass der Massenlauf $m(\theta)$ auf dem
Kreis stetig schließt. Die einfachste Schließung ist ein
\emph{gerades} Profil um den Antipoden: Minimum bei $\theta=\pi$,
symmetrische Rückkehr auf der anderen Hälfte. Glatt ist sie, wenn
$dm/d\theta \to 0$ am Umkehrpunkt.

In der Strahlungsära gilt $H \simeq H_0\sqrt{\Omega_r}(1{+}z)^2$,
daraus analytisch

\[
  \eta_0 - \chi = \frac{c}{H_0\sqrt{\Omega_r}}\,\frac{1}{1{+}z}
  \qquad\Longrightarrow\qquad
  \left|\frac{dm}{d\theta}\right|_{\text{Ende}}
  = \sqrt{\Omega_r} \approx 0{,}0096.
\]

\textbf{Befund:} Das beobachtete Profil läuft mit Steigung
$\sim 0{,}01$ in den Endpunkt -- die Glattheitsbedingung einer
geraden periodischen Schließung ist auf $1\,\%$ erfüllt. Die
nötige Deformation gegenüber der gemessenen Geschichte ist ein
$1\,\%$-Effekt am äußersten Rand.

\section{Warum das kein Selbstbetrug ist}

Die Rückkehr-Hälfte -- dort, wo der Massenlauf wieder ansteigen
müsste -- ist \emph{exakt} die Hälfte hinter der Streuwand:
photonisch prinzipiell unbeobachtbar. Das ist zunächst verdächtig
bequem, aber es ist keine Wahl dieses Dokuments, sondern Folge des
Antipoden-Befunds: Die Wand \emph{liegt} bei $0{,}992\,\pi$. Die
Schließung widerspricht deshalb keiner einzigen Photon-Beobachtung
-- und sie ist trotzdem nicht unprüfbar (Kap.~E).

Das Klimazonen-Bild aus der Diskussion zu Dok.~312 wird damit
quantitativ: Der Gradient ist die Landkarte \emph{entlang} des
Zyklus (wie Klimazonen entlang eines Meridiankreises), keine
Entwicklung \emph{aus} einem Anfang -- und der \enquote{heißeste}
Punkt der Karte ist der Antipode.

% ============================================================
\chapter{Das Entropie-Hindernis: der Kern von D(ii) [K]}
% ============================================================

\section{Periodizität ist nicht generisch}

Das Entropiebudget des beobachtbaren Volumens ist
$S \sim 10^{104}\,k_B$ (dominiert von supermassereichen Schwarzen
Löchern; Photonen: $10^{89}\,k_B$). Poincaré-Wiederkehr eines
generischen Zustands braucht $\sim e^{S/k_B} \sim 10^{10^{104}}$
dynamische Zeiten -- gegen $\tau_c \sim 10^{18{,}5}$~s. Die
Diskrepanz beträgt $\sim 10^{104}$ \emph{Größenordnungen im
Exponenten}: Exakte $\tau_c$-Periodizität kann dynamisch nicht
erwartet werden.

\textbf{Konsequenz:} Periodizität ist eine
\textbf{Randbedingung} -- eine Auswahl periodischer Lösungen aus
dem Lösungsraum, keine Folge der Evolution. Das ist dieselbe
Kategorie wie die Wahl der T$^4$-Topologie selbst ([SETZUNG] im
Sinne von Dok.~311) und muss so gebucht werden.

\section{Fein- und grobkörnig getrennt}

\textbf{Feinkörnig [K]:} Unter unitärer Evolution ist die
feinkörnige Entropie konstant; einer exakt periodischen Lösung
steht mikroskopisch nichts entgegen. Die Randbedingung ist
zulässig.

\textbf{Grobkörnig [offen]:} Der Zweite Hauptsatz ist eine
Aussage über grobkörnige Entropie relativ zu einem Beobachter.
Auf der Rückkehr-Hälfte muss sie abnehmen -- was beobachterrelativ
möglich ist (Umkehr der Korrelationsrichtung; keiner
\emph{dort} sähe eine Verletzung, da der thermodynamische
Zeitpfeil mitdreht), aber unbewiesen. \textbf{Genau das ist der
präzisierte Kern von D(ii):} nicht die Feldperiodizität (die ist
als Randbedingung konsistent, Kap.~C), sondern der Nachweis, dass
eine grobkörnige Entropie-Schließung mit der beobachteten
Strukturbildung verträglich ist. Bis dahin: offen, ohne
Beschönigung.

% ============================================================
\chapter{Prüfflächen [S]}
% ============================================================

Photonen enden an der Wand ($0{,}992\,\pi$). Zwei Boten stammen
vom Antipoden selbst oder dahinter:

\begin{center}
\begin{tabular}{lll}
\toprule
Bote & Entkopplung & Position \\
\midrule
C$\nu$B (Neutrino-Hintergrund) & $z \sim 6\times10^{9}$ &
  $\theta \approx 1{,}01\,\pi$ \\
primordiale GW & nominell $z \sim 10^{25}$ &
  $\theta \approx 1{,}01\,\pi$ \\
\bottomrule
\end{tabular}
\end{center}

Eine \emph{nicht-monotone} Signatur dort -- die Umkehr des
Massenlaufs jenseits des Antipoden -- wäre die prinzipielle
Prüffläche der Schließung. Praktisch außerhalb heutiger
Messbarkeit, aber benennbar; zusammen mit dem
$4{,}2'$-Dipolversatz (Dok.~312, Torus-Ruhesystem) sind das zwei
konkrete Signaturen der kompakten Lesart.

% ============================================================
\chapter{Statusübersicht}
% ============================================================

\begin{center}
\begin{tabular}{lp{6.8cm}l}
\toprule
Aussage & Inhalt & Status \\
\midrule
Antipoden-Befund & Geschichte $=$ ein Halbzyklus auf $1{,}2\,\%$;
  Streuwand bei $0{,}992\,\pi$ & [K$|$P20] \\
Koinzidenz-Einordnung & $\eta_0 H_0/c \approx \pi$: Strukturaussage
  statt Zufall erst nach Vorwärtsherleitung & [S] \\
Glatte Schließung & Endpunktsteigung $\sqrt{\Omega_r}
  \approx 0{,}01$; gerade Schließung auf $1\,\%$ erfüllt & [K] \\
Rückkehr-Hälfte & $=$ photonisch unbeobachtbare Hälfte;
  kein Konflikt mit Photon-Daten & [K] \\
Periodizität & Randbedingung (Poincaré-Diskrepanz
  $10^{104}$~dex im Exponenten), nicht Dynamik & [K] \\
Entropie & feinkörnig zulässig; grobkörnige Schließung offen --
  präzisierter Kern von D(ii) & offen \\
Prüfflächen & C$\nu$B/GW vom Antipoden; Nicht-Monotonie als
  Signatur & [S] \\
\bottomrule
\end{tabular}
\end{center}

\medskip
\noindent\textbf{Bilanz:} \enquote{Kein Urknall} heißt nach diesem
Dokument präzise: Der \enquote{Anfang} ist der Antipode des
Zeitzyklus -- ein Ort, kein Ereignis; die beobachtete Geschichte
ist die eine Hälfte einer auf $1\,\%$ glatt schließbaren Landkarte;
und die verbleibende Last liegt nicht bei den Feldern, sondern bei
der grobkörnigen Entropie. D(ii) ist damit nicht erledigt, sondern
auf seinen harten Kern reduziert.

\medskip
\noindent\textbf{Registrierung:} Aufnahme in Dok.~190 erfolgt
separat.

\end{document}
